\subsection{Bei größerer Tiefe reichen Vorzeichen nicht aus}
Für eine ganzzahlige Inklusion $M=K\Z^8\subseteq L$ und einen Kindcharakter $c$ ergibt
\begin{equation}
\theta=K^{-T}c,\qquad
\chi_\theta(x)=e^{\pi i\theta\cdot x},\qquad
\chi_\theta(Kz)=(-1)^{c\cdot z}
\end{equation}
eine Fortsetzung auf derselben Quelle. Dabei können Phasenordnungen $2,2,4,4,8$ auftreten. Für Ordnung vier ist $(I+U)/2$ kein Projektor. Der richtige Sektorprojektor lautet bei Ordnung $N$
\begin{equation}
\Pi_r=\frac1N\sum_{k=0}^{N-1}e^{-2\pi i rk/N}U_\theta^k.
\end{equation}
Die kohärente Speicherung
\begin{equation}
V\psi=\sum_r\ket r\otimes\Pi_r\psi
\end{equation}
ist isometrisch. Für $\theta/2=\ell/N$ mit $\ell\in\Z^8$ ist das Sektorlabel $r(x)=\ell\cdot x\pmod N$. Wenn ein Feld die Ladung um $\alpha$ ändert, muss sein Verbraucher den Speicher gleichzeitig verschieben:
\begin{equation}
\widetilde E_\alpha=S_{r(\alpha)}\otimes E_\alpha,
\qquad\widetilde E_\alpha V=VE_\alpha.
\end{equation}
Ein gespeichertes Label ohne diesen Verbraucher wäre keine vollständige Rekursion. Wird der Speicher verworfen, entsteht die Dephasierung $\rho\mapsto\sum_r\Pi_r\rho\Pi_r$; die Information ist dann tatsächlich verloren.

\begin{center}
\begin{tikzpicture}[node distance=9mm]
\node[box,text width=3.1cm] (s) {voller Zustand\\Ladung + Phase};
\node[box,right=of s,text width=4.1cm] (c) {reduzierte Darstellung\\plus kohärenter Carry};
\node[box,right=of c,text width=3.1cm] (e) {geladene Operation\\ändert beide Teile};
\draw[arr] (s)--(c);\draw[arr] (c)--(e);
\node[below=7mm of c,align=center,font=\small,text width=13cm] {Carry ist die mitgeführte Information, auf die ein späterer Test reagieren kann.\\Er ist kein kostenloser Ersatz für Gedächtnis.};
\end{tikzpicture}
\end{center}

Nach $n$ echten Index-zwei-Inklusionen mit der geerbten Gitterform ist die arithmetische Bilanz
\begin{equation}
[L:M_n]=2^n,\qquad |M_n^\#/M_n|=4^n,\qquad |L/M_n|=2^n.
\end{equation}
Die ursprüngliche Quelle belegt nur einen bestimmten Teil der gesamten dualen Klassen. Ihre geladene Vervollständigung rekonstruiert die Quelle. Der Index $2^n$ ist keine abgeleitete Anzahl räumlicher Zellen.

\subsection{Der Lift muss die wirkliche Operatorwirkung erhalten}
Auf dem unverdrehten Cartanraum wird die native Reflexion durch $s_\alpha s_{J\alpha}$ realisiert. Die gleiche Pauli-/Heisenberg-Paarung eines verdrehten Moduls reicht dagegen nicht: Die native Reflexion hat eine $3+1$-Aufteilung mit $|\tr|^2=4$, während die untersuchten Pauli-ambigen Cliffordimplementer $2+2$ oder $1+1+1+1$ und $|\tr|^2=8$ oder $0$ besitzen. Das schließt die behauptete Identifikation dieser Wirkungen aus.

Geladene Cartanlifts benötigen Charakterphasen. Die Wörter $0,1,0$ und $2,1,2$ haben dieselbe nackte Gitterwirkung, unterscheiden sich aber um Deckcharakter 152. Eine fortsetzbare Darstellung der geprüften endlichen Cartan-Liftgruppe speichert deshalb Gitterwirkung plus acht Charakterbits. Die allgemeinen Hecke-Fortsetzungen benötigen dagegen rationale Phasenkoordinaten $\theta\bmod2$. Sie braucht nicht die gesamte alte Wortliste. Ein ganzzahliges Relationszeugnis schließt in der geprüften Liftaufgabe eine globale Sektion auch unter kontinuierlichen Toruskorrekturen aus.

Der ältere Quartiklift ist ein anderer tatsächlicher $E_8$-Automorphismus:
\begin{equation}
g_5=-T_e,\qquad g_4=e^{-i\pi/4}r_e,
\qquad U_X=\overline g_5\otimes g_4.
\end{equation}
Auf $X$ gilt $U_X^2=-iI$; der Lift hat Ordnung acht. Seine adjungierte Spur $16$ unterscheidet ihn von den involutiven Cartanlifts mit Spuren $-8$ oder $24$. Nur 16 der 60 nativen Transporte erhalten den fest gewählten P2-Generator; die übrigen transformieren den Ladungsrahmen kovariant. Sie sind deshalb nicht automatisch physische Ereignisse, die dieselbe feste Hyperladung erhalten.

\section{Universalraum: Der Raum dessen, was eine Zukunft noch unterscheiden kann}\label{sec:universal}

\key{Zwei Vorgeschichten gehören genau dann zum selben wirksamen Zustand, wenn keine zugelassene spätere Handlung oder Messung sie unterscheiden kann. Welcher Universalraum gemeint ist, hängt deshalb vom tatsächlichen Operations- und Auslesesystem ab.}

\subsection{Operationale Äquivalenz und lineare Vorhersage}
Für eine festgelegte Klasse zulässiger Zukunftstests $T$ ist
\begin{equation}
h\sim h'\quad\Longleftrightarrow\quad
\mathcal P(T\mid h)=\mathcal P(T\mid h')\quad\text{für alle }T.
\end{equation}
Wenn das Gewicht der vorbereitenden Geschichte relevant ist, müssen unnormierte gemeinsame Antworten verglichen werden. Wenn kohärente Referenzen und Interferenz zulässig sind, gehören auch diese zum Testsystem. Eine Aussage über alle Zukunftstests darf nicht auf die gerade sichtbare Bildschirmablesung verkürzt werden.

Für eine skalare Wortantwort $s$ bildet
\begin{equation}
H_{u,v}=s(uv)
\end{equation}
die Hankelmatrix. Ihre Zeilen sind Zukunftsantworten von Vorgeschichten. Bei endlichem Rang entspricht $\rank H$ der Dimension einer minimalen linearen Wortrealisierung unter den üblichen algebraischen Voraussetzungen \cite{carlyle,thon}. Daraus folgen weder Positivität noch automatisch eine minimale Quantensystemdimension.

\subsection{Die positive kohärente Rekonstruktion}
Für Operatorwörter $K_u$ ist der kohärente Kern
\begin{equation}
D(u,v)=\omega(K_u^\dagger K_v).
\end{equation}
Er enthält Kreuzterme zwischen Geschichten. Auf einer unitalen $C^*$-Algebra mit normiertem positivem Zustand liefert die GNS-Konstruktion
\begin{equation}
\mathcal N_\omega=\{A:\omega(A^\dagger A)=0\},\qquad
\mathcal H_\omega=\overline{\mathfrak A/\mathcal N_\omega}.
\end{equation}
Das Nullideal ist unter Linksmultiplikation stabil; die zyklische Darstellung ist bis auf unitäre Äquivalenz eindeutig \cite{segal}. Bei einer allgemeinen Algebra unbeschränkter Felder sind gemeinsame Domänen und analytische Voraussetzungen zusätzlich zu behandeln. Die pauschale Behauptung einer Darstellung in beschränkten Operatoren für jede positive Feld-$*$-Algebra wäre zu stark.

Die GNS-Konstruktion rekonstruiert eine Darstellung aus einer gegebenen positiven Quelle. Sie wählt nicht voraussetzungslos die Quelle, ihre Positivität, die Born-Auslese oder die physisch zulässigen Eingriffe. Ein beliebiger positiver Gramkern auf Geschichtslabels ist außerdem noch kein kausal normalisierter Mehrzeitprozess. Prozess-Tensoren beziehungsweise Quantennetzwerke verlangen zusätzliche Kompositions- und Normierungsbedingungen \cite{combs}.

\subsection{Ein Beispiel, das die drei Ebenen sofort trennt}
Das implementierte Instrument $K_e=U_e/\sqrt{15}$ ergibt bei alleiniger Beobachtung seiner Ereignislabels
\begin{equation}
p(e_n\cdots e_1)=15^{-n}.
\end{equation}
Die zugehörige klassische Wort-Hankelmatrix hat Rang eins. Für denselben Operatorprozess mit $\omega(A)=\tr(A)/5$ besitzt der kohärente Kern dagegen den berechneten GNS-Rang 25, weil die Operatorwörter $M_5(\C)$ aufspannen.

Das ist kein Widerspruch. Die Labelauslese sieht eine einfache Folge gleichgewichteter Ereignisse. Ein kohärenter Zukunftstest kann zusätzliche Operatorunterschiede erkennen. Ebenso erzeugen $I$ und $-I$ dieselbe isolierte Quantenkanalkarte, aber verschiedene GNS-Vektoren; erst eine zugelassene kohärente Referenz macht eine relative Phase sichtbar.

Damit wird die stärkste eingefügte Universalprozess-Synthese präzisiert: Operationale Zukunftsklassen, lineare Hankelrealisierung und positiver GNS-Raum sind drei verbundene Konstruktionen, deren Gleichsetzung passende Tests und Voraussetzungen benötigt.

\subsection{Die exakte Bedingung für eine verlustfreie Rekodierung}
Eine Reduktion $R$ ist unter einer Fortsetzung $F$ geschlossen, wenn
\begin{equation}\label{eq:quotient}
RF=\overline F R.
\end{equation}
Im linearen Fall ist insbesondere $F(\ker R)\subseteq\ker R$ erforderlich. Zusätzlich müssen alle relevanten Auslesungen durch $R$ faktorisieren. Dann folgt durch Induktion die Gleichheit der Antworten für beliebig lange zulässige Operationswörter.

Die Hecke-Carry-Intertwiner liefern eine konkrete Instanz dieser Regel. Der Universalraumanschluss ist deshalb nicht bloß eine Metapher. Er sagt genau, welche Information auf einer gröberen Darstellung bleiben und wie sie bei einer Operation weiterwirken muss.

\subsection{Warum ein Fixpunkt des Gesetzes keine universelle Vergessensmaschine ist}
Ein primitiver positiver Operator kann unter passenden endlichen Voraussetzungen einen eindeutigen normierten Fixpunkt besitzen. Doch P1-Primitivität ist nicht automatisch diese Operatorprimitivität; auch ein endlicher Satz darf nicht ohne Grenzwertkontrolle auf einen unendlichen Raum übertragen werden.

Für einen CPTP-Kanal $\mathcal R$ mit CPTP-Linksinverser $\mathcal L$ auf der gesamten betrachteten Zustandsklasse gilt
\begin{equation}
\mathcal L\mathcal R=\Id
\quad\Longrightarrow\quad
\norm{\mathcal R(\rho)-\mathcal R(\sigma)}_1
=\norm{\rho-\sigma}_1.
\end{equation}
Beide Ungleichungen folgen aus der Kontraktivität der Spurweite. Eine derart verlustfreie Abbildung auf einem nichttrivialen festen Zustandsraum kann deshalb nicht zugleich alle unterscheidbaren Zustände in denselben primitiven Attraktor treiben.

Ein eindeutiges \emph{Gesetz} kann viele Zustände haben. Ein reduzierter Kanal kann nach bewusstem Vergessen relaxieren, während die volle Entwicklung einschließlich Umgebung Information bewahrt. Zukunftsvollständigkeit des Prozessraums allein behauptet keine Linksinvertierbarkeit seiner Vergröberung. Diese Unterscheidungen verhindern einen falschen Kurzschluss von ``rekursiv konsistent'' zu ``einziger physischer Weltzustand''.

\section{Raum, Topologie und Ausbreitung: Was aus dem Graphen folgt}

\subsection{Die räumliche Konstruktion ist markiert}
Die Tutte--Coxeter-Zelle besitzt dreißig Knotentypen, fünfundvierzig Kanten, Grad drei und Zyklusrang sechzehn. Der verwendete markierte periodische $A_3$-Cover wählt daraus einen Rang-drei-Periodenquotienten. In der benutzten Normierung ergibt die kleine Blochwelle
\begin{equation}
\lambda_{\min}(k)=\frac85\norm{k}^2+O(\norm{k}^4).
\end{equation}
Dies ist ein konkreter dreidimensionaler Ausbreitungszweig. Die Dimension stammt aus der gewählten Periodenmarkierung, weder aus den sechzehn unabhängigen Graphschleifen noch aus dem Rang acht des internen Ladungsgitters. Die universelle Überlagerung des endlichen Graphen wäre dagegen ein Baum.

Zehn disjunkte markierte Pfade decken die dreißig Typen. Fünf Blocktypen tragen $U$, fünf $\overline U$. Damit besitzt die ladungskovariante Pfadrekursion einen wirklichen Ort in diesem markierten Netz. Die Auswahl der Perioden und ihre physische Bedeutung bleiben als Voraussetzungen sichtbar.

\subsection{Die effektive Metrikantwort braucht die gesamte Zelle}
Ein naiver Ausdruck aus lokalen zweiten Sprungmomenten,
$\sum_e r_e\,\Delta x_e\Delta x_e^T$, muss noch keine globale effektive Metrik ergeben. An einzelnen Knoten ist er sogar singulär. Der interne harmonische Korrektor beschreibt, wie sich die Knotentypen einer ganzen Periode gemeinsam einstellen.

Die überprüfbare effektive Antwort wird deshalb aus der Blochkrümmung gewonnen:
\begin{equation}
D_{ij}=\frac12\left.\frac{\partial^2\lambda_{\min}(k)}{\partial k_i\partial k_j}\right|_{k=0}.
\end{equation}
In der isotropen Ausgangsnormierung ist $D=(8/5)I_3$. Die Sitzung berechnet auch die Antwort auf geänderte Kantenraten mit dem vollständigen Zellkorrektor. Das ist mehr als eine gezeichnete Raumvermutung, aber noch keine abgeleitete Lorentzmetrik. Diffusive beziehungsweise spektrale Kleinimpulsantwort, unitäre Wellenausbreitung und relativistische Kausalstruktur müssen passend unterschieden werden.

\picturetext{Die scheinbare Weglänge in einem Straßennetz hängt nicht nur von einer Kreuzung ab. Manche Wege zwingen zu Umwegen innerhalb des Viertels. Erst wenn das ganze wiederkehrende Viertel berücksichtigt ist, erhält man die effektive Ausbreitung über große Distanzen.}

\subsection{Existenzsätze erhalten, ihre Reichweite nicht ausweiten}
Die ältere thermodynamische Dynamikarbeit enthält bereits Lieb--Robinson-Schranken, Zeitentwicklung, Gauss-Kompatibilität sowie Grundzustands-/KMS-Existenz für die dort spezifizierte lokale Hamiltonklasse mit endlichen Faktoren auf $\Z^2$. Diese Ergebnisse werden hier nicht erneut als ungelöste reine Existenzfragen behandelt. Sie begründen aber auch nicht von selbst die ausgewählte räumliche $3+1$-Realisierung der vollständigen geladenen Quelle.

Die richtige Anschlussfrage lautet: Welcher dieser bereits verstandenen lokalen Mechanismen trägt genau die jetzt festgelegte Quelle, ihre ursprünglichen Markierungen und ihre physischen Auslesungen? Dabei müssen dieselbe Zeit und derselbe Zustand durch die Konstruktion laufen.

\section{Zuse und Wolfram: Der produktive gemeinsame Kern und die nötigen Unterschiede}

Zuses \emph{Rechnender Raum} motiviert den Gedanken, physische Entwicklung als lokale regelhafte Verarbeitung zu beschreiben. Für diese Sitzung ist die nützliche Konsequenz: Eine Theorie muss sagen, was einen Zustand ausmacht, welcher lokale Schritt erlaubt ist und wie große Strukturen aus wiederholter Komposition entstehen \cite{zuse}. Die konkrete TFPT-Quelle liefert dafür eine reichere Algebra als ein bloßes Bitgitter; daraus folgt keine historische Prioritätsbehauptung und kein Beweis, dass das Universum ein bestimmter Zellularautomat ist.

Wolframs Multiway- und Hypergraphbilder helfen, Zustände, mögliche Aktualisierungen und Abhängigkeiten getrennt darzustellen. Seine eigene Behandlung gewichteter Multiwaygraphen zeigt bereits, dass eine gleichmäßige Verteilung auf Nachfolger bei anderer Foliation keine robuste Normierung garantieren muss \cite{wolfram}. Deshalb darf eine graphische Verzweigung nicht stillschweigend ihre physikalischen Wahrscheinlichkeiten festlegen.

Die konkrete Verbindung zu den Ergebnissen dieser Sitzung ist die folgende Typentrennung:
\begin{longtable}{L{0.21\linewidth}L{0.34\linewidth}L{0.37\linewidth}}
\toprule
Graphart & Knoten und Kanten & Was zusätzlich zu prüfen ist\\\midrule\endhead
Ladungsgraph & Wurzeln, Cartanrichtungen, nichtverschwindende Klammern & Wirkliche Darstellung, Kokzyklus, Markierungen\\
Rechengraph & Teilwörter und Ward-Reduktionsschritte & Äquivalenz verschiedener Auswertungsbäume\\
Antwort-Hypergraph & Verbundene Mehrfeldkoeffizienten & Quelle, Zustand, Positionen; keine automatische Kausalinterpretation\\
Geschichtsgraph & Physische Operationen und ihre Verkettung & Instrumente, Phasen, kausale Normierung\\
Raumgraph & Lokale Zellen und räumliche Nachbarschaft & Periodenwahl, Ausbreitung, Kontinuums- und Zeitanschluss\\
Theoriegraph & Aussagen und ihre Begründungsabhängigkeiten & Originalbeleg und Voraussetzungen; ein Pfad ist kein Beweis\\\bottomrule
\end{longtable}

Die besonders wichtige positive Verbindung ist: Das ausgeführte Quellgesetz besitzt unterschiedliche exakte Auswertungsbäume, während sein verbundener Anteil echte höhere Antwort-Hyperkanten liefert. Beides stammt aus derselben Algebra. Die Bäume dürfen optimiert werden; die Quellantwort muss erhalten bleiben. Werden diese Auswertungsbäume stattdessen als verschiedene Welthistorien gezählt und erneut gewichtet, ändert man das bereits festgelegte Gesetz.

Eine gute Gesamtkonstruktion würde die Typen durch explizite Abbildungen verbinden: Ein räumliches lokales Ereignis wirkt durch einen Operator; dieser verändert einen Zustand und seinen Carry; das erzeugt eine überprüfbare Mehrzeitantwort; eine gröbere Darstellung berechnet dieselbe Antwort. Der Graph ist dann eine Darstellung der vollständigen Regel und kein Ersatz für sie.

\section{Fortsetzung: Die Auswahl von Level eins ist auf einen Ursprungssatz reduziert}\label{sec:levelselection}

\key{\textbf{Ein positiver Auswahlsatz ist jetzt verfügbar.} Innerhalb einer unitären positiven affinen $E_8$-Quelle genügt entweder eine unabhängig begründete Zentralzahl höchstens acht oder die Nullrelation eines einzigen echten langen Wurzelstroms, um Level eins zu erzwingen. Das verbleibende Problem ist die Herkunft dieser konkreten Eigenschaft aus der ursprünglichen Quelle.}

\subsection{Was Zentralzahl acht tatsächlich erzwingt}
Die Vergleichsklasse ist ausdrücklich festgelegt: eine einfache unitäre Vertexoperatoralgebra vom CFT-Typ, mit positivem Energiegrad, eindeutigem Vakuum im Grad null und einer tatsächlichen kompakten $E_8$-Stromalgebra im Grad eins. Lange Wurzeln haben Normquadrat zwei. Der affine Level ist dann eine positive ganze Zahl $k$; die von den Strömen erzeugte Algebra ist $L_k(\mathfrak e_8)$ \cite{donglin}.

Ihr Sugawara-Stress-Tensor hat
\begin{equation}
c_{\rm aff}(k)=\frac{248k}{k+30},\qquad
c_{\rm ges}\ge c_{\rm aff}(k)\ge8.
\end{equation}
Die erste Ungleichung folgt aus der Positivität der verbleibenden konformen Komponente. Wenn der Ursprung unabhängig $c_{\rm ges}\le8$ liefert, folgt zwingend $k=1$ und $c_{\rm ges}=8$. Noch stärker:
\begin{equation}
\norm{\omega_{\rm ges}-\omega_{\rm Sug}}^2
=\frac{c_{\rm ges}-c_{\rm aff}}2=0.
\end{equation}
Damit stimmen die Stress-Tensoren selbst überein. Bei Erzeugung durch die Ströme ist die Quelle $L_1(\mathfrak e_8)\cong V_{E_8}$; der Übergang zum lokalen Netz benötigt die einschlägigen starken Lokalitätsvoraussetzungen \cite{cklw}. Die Eindeutigkeit des affinen Vakuumsektors wählt noch keinen angeregten physischen oder kosmologischen Zustand aus.

\textbf{Die kritische Richtung:} Rang acht gibt zunächst eine Untergrenze für die Zentralzahl, keine Obergrenze. Auch ``nur diese Ströme, keine entkoppelten Zusatzsysteme'' genügt nicht: $L_2(\mathfrak e_8)$ ist eine unitäre, von denselben 248 Stromrichtungen erzeugte Alternative mit $c=31/2$. Sie ist ein Gegenfall zu genau diesem schwächeren Auswahlargument, keine Behauptung, sämtliche übrigen TFPT-Ursprungsbedingungen zu erfüllen.

Die kanonische Gitterquantisierung aus Abschnitt~\ref{sec:source} liefert ihrerseits bereits Level eins: Der Wurzelvertex besitzt bei Normquadrat zwei einen Doppelpol. Sobald diese Quantisierung als ursprüngliche Quelle identifiziert ist, ist ihr Level keine weitere freie Wahl. Genau diese Identifikation darf aber nicht durch den schon daraus berechneten Doppelpol bewiesen werden.

\subsection{Ein einziger tatsächlicher Wurzelstrom genügt}
Für einen normierten langen Wurzelstrom und sein wirkliches Adjungiertes gilt
\begin{equation}
[e_m,f_n]=h_{m+n}+mk\delta_{m+n,0},\quad
[h_0,e_{-1}]=2e_{-1},\quad e_{-1}^{\dagger}=f_1.
\end{equation}
Anwendung auf das Vakuum ergibt
\begin{equation}\label{eq:rootnull}
\norm{e_{-1}^2\Omega}^2=2k(k-1).
\end{equation}
Somit erzwingen $e_{-1}\Omega\ne0$ und $e_{-1}^2\Omega=0$ den Level eins. Die zweite Gleichung ist eine Aussage über einen vollständigen affinen Modus. Die Nilpotenz irgendeiner endlichen Pauli-Matrix würde dafür nicht genügen.

Diese Nullrelation besitzt eine konkrete ursprüngliche Realisierung, sobald zwei unabhängige gleichchirale CAR-Arten vorhanden sind. Mit der tatsächlichen Viertelholonomie-Polarisation
\begin{equation}
a_{i,j}\Omega=0\ (j\le0),\qquad
a^\dagger_{i,j}\Omega=0\ (j\ge1)
\end{equation}
liefert der neutrale Wurzelstrom
\begin{equation}
E_{-1}=\sum_j a^\dagger_{1,j}a_{2,j+1},\qquad
E_{-1}\Omega=a^\dagger_{1,0}a_{2,1}\Omega,\qquad E_{-1}^2\Omega=0.
\end{equation}
Seine erste Norm ist eins; seine Energie ist $1/4+3/4=1$. Die gemeinsame chemische Viertelverschiebung fällt im neutralen Strom weg.

Der bereits hergeleitete einzelne Rohkanal liefert dagegen einen Dichtestrom $J_{-1}=\sum_j a_j^\dagger a_{j+1}$ mit
\begin{equation}
J_{-1}^2\Omega=a_{-1}^\dagger a_1\Omega-a_0^\dagger a_2\Omega,
\qquad\norm{J_{-1}^2\Omega}^2=2.
\end{equation}
Beide Rechnungen wurden mit den ursprünglichen CAR-Matrizen geprüft; das Fenster $j=-1,0,1,2$ enthält für diese zwei Anwendungen alle nichtverschwindenden Beiträge. Ein einzelner Kanal darf deshalb nicht als das benötigte unabhängige Artenpaar ausgegeben werden. Der vorhandene Kanalvertrag findet einen komplexen chiralen Kanal je Rand, mit entgegengesetzter Orientierung der Ränder; die acht Reihen sind keine acht unabhängigen gleichchiralen Ströme.

\picturetext{Die ganze Anlage muss nicht nochmals aus 248 Einzelteilen erraten werden. Ein einziges richtig angeschlossenes Zahnrad kann die gemeinsame Übersetzung festlegen. Es muss jedoch wirklich ein Zahnrad derselben Anlage sein. Ein ähnlich aussehender Schalter erfüllt die benötigte Relation nicht.}

Der kleinste hinreichende Herkunftssatz lautet damit: Die Rohquelle liefert einen vakuum-, moden-, adjungierten- und ladungsverträglichen langen Wurzelstrom mit Einbettungsindex eins und der Nullrelation~\eqref{eq:rootnull}. Die alte freie R2-Route mit 16 reellen Majoranas liefert eine hinreichende Mehrkanalrealisierung, sobald ihre Quellidentifikation gilt. Ihre Herkunft wird durch die bloße Zahl 16 nicht nachgeholt.

\subsection{Warum die bisherige Fünferrekursion den Level nicht auswählt}
Die tatsächliche Chevalley-Ward-Rechnung lässt sich ohne Level-eins-Vorentscheidung ausführen. Für die Felder $X_{ia},X_{jb}^\dagger,X_{kc},X_{ld}^\dagger$ an $(-1,0,1,\infty)$ folgt exakt
\begin{equation}
F_k=k^2(A+D)+k(B+C),
\end{equation}
mit
\begin{align}
A&=\delta_{ij}\delta_{kl}\delta_{ab}\delta_{cd},&
D&=\delta_{il}\delta_{jk}\delta_{ad}\delta_{bc},\\
B&=\delta_{ij}\delta_{kl}\delta_{ad}\delta_{bc},&
C&=\delta_{il}\delta_{jk}\delta_{ab}\delta_{cd}.
\end{align}
Bei festem Familienindex bleibt $k(k+1)(\delta_{ij}\delta_{kl}+\delta_{il}\delta_{jk})$. Nach Gesamtnormierung ist der Fünfertensor daher für alle positiven $k$ derselbe. Die normierte Fünferrekursion bestätigt ihre Form, sieht aber diese Levelalternative nicht.

Zwei gemischte Antworten unterscheiden sie unmittelbar: Für Carrierindizes $(0,0,1,1)$ ergeben Familienindizes $(0,0,1,1)$ den Wert $k^2$ und $(0,1,1,0)$ den Wert $k$. Das Verhältnis ist $1/k$. Es bleibt unter einer gemeinsamen Feldnormierung erhalten.

\begin{center}
\begin{tabular}{ccccc}
\toprule
$k$ & $c_{\rm aff}$ & gemischtes Verhältnis & $\norm{e_{-1}^2\Omega}^2$ & neutrale Restnorm$^2$\\\midrule
1 & 8 & 1 & 0 & 0\\
2 & $31/2$ & $1/2$ & 4 & $15/4$\\
3 & $248/11$ & $1/3$ & 12 & $80/11$\\\bottomrule
\end{tabular}
\end{center}
Die letzte Spalte ist $\norm{\omega_{\rm Sug}-\omega_{\rm Heis}}^2=120(k-1)/(k+30)$. Im Grad zwei besitzt die Level-eins-Gitterquelle 4124 Zustände, die affinen Quellen mit $k\ge2$ dagegen 31124. Die Differenz 27000 entsteht durch die Level-eins-Nullrelationen. Diese exakten Kennzeichen ersetzen keine unabhängige Herleitung ihrer Erfüllung durch den Rohprozess.

Die neue Rechnung prüft alle 225 Paare von Gleichheitsmustern der vier Carrier- und Familienindizes; die Erweiterung auf konkrete Indexbelegungen verwendet die bereits geprüfte Tensor-Kovarianz. Das volle $W_5\otimes W_4$ fordert im affinen Vergleich tatsächlich $k=1$. Wurde dieser Tensor zuvor aus $E_{8,1}$ hergeleitet, ist er aber kein unabhängiger Ursprungsbeweis für dieselbe Quelle.

\subsection{Primitive Windung und Determinanten haben unterschiedliche Normierungen}
Die alte Rechnung \file{v470_alpha_inflow_level.py} identifiziert $|C|=1$ mit dem primitiven $U(1)$-Inflow und dem betreffenden Koeffizienten der $\alpha$-Gleichung. Sie identifiziert nicht ohne zusätzlichen Abbildungssatz den affinen $E_8$-Level.

Bei einer Determinante in einer endlichen Darstellung $\rho$ multipliziert der Dynkinindex $d_\rho$ den grundlegenden affinen Kokzyklus. In der Normierung langer Wurzeln auf zwei hat die adjungierte 248er-Darstellung von $E_8$ Index 60; der kleinste nichtverschwindende Darstellungsindex ist 60 \cite{laszlosorger}. Deshalb ist ``eine adjungierte Fermiondeterminante'' kein automatischer Level-eins-Beweis. Die ursprüngliche Determinantenlinie muss auf die grundlegende Linie samt Verbindung mit der tatsächlich passenden Normierung abgebildet werden. Die bereits formulierte kontinuierliche Quillen-Brücke bleibt genau dafür relevant.

\section{Konstruktiver Abschluss innerhalb der Quelle: Felder, Antworten und Gedächtnis}\label{sec:constructive}

\key{\textbf{Die Quelle hat jetzt eine zweite vollständige, ausführbare Darstellung.} Die $X$-Felder sind explizite Fermionbilineare; die $C$-Felder sind die erforderlichen geladenen Erweiterungsfelder. Mit dem ursprünglichen Kokzyklus ergibt sich eine geschlossene Formel für sämtliche endlichen C/X-Wörter. Dieselbe Formel liefert einen positiven Geschichtskern und eine exakte Darstellung ihres gesamten Nachfahrengedächtnisses.}

\subsection{Die neun Hilfsfermionen und die vier Erweiterungsfelder}
Die aus den ursprünglichen Wurzeln gebaute $A_8$-Isometrie bildet alle 20 X-Wurzeln exakt auf
\begin{equation}
\alpha(X_{ia})=e_{5+a}-e_i\quad(0\le i<5,\ 0\le a<4)
\end{equation}
im neundimensionalen Nullsummenraum ab. Neun komplexe chirale Hilfsfermionen realisieren deshalb die $A_{8,1}$-Ströme durch
\begin{equation}
X_{ia}(z)=:\!\psi_{5+a}^\dagger(z)\psi_i(z)\!:.
\end{equation}
Der gemeinsame $U(1)$-Anteil wird abgespalten: $c=9-1=8$. Die tatsächlichen 8000 Tripelklammern erfüllen
\begin{equation}
[[X_{ia},X_{jb}^\dagger],X_{kc}]
=\delta_{ij}\delta_{bc}X_{ka}+\delta_{ab}\delta_{jk}X_{ic}.
\end{equation}
Die ursprünglichen Phasen passen also zu den Matrixeinheiten; es genügt nicht nur eine Dimensionsgleichheit. Die P2-Wirkung wird mitgeführt: Die diagonalen Werte
\[
(-1/3,-1/3,-1/3,1/2,1/2,0,0,0,0)
\]
geben auf $X_{ia}$ genau $-y_i$.

Die vier C-Felder besitzen in denselben Koordinaten die Ladungen
\begin{equation}\label{eq:Csu9}
q(C_a)=\sum_{i=0}^4e_i+e_{5+a}-\frac23\sum_{r=0}^8e_r.
\end{equation}
Sechs Einträge sind $1/3$, drei sind $-2/3$; die Norm ist zwei und die P2-Ladung null. Das sind Gewichte des $\Lambda_6$-Sektors. Das zugehörige Sechsfermionprodukt hätte zunächst Gewicht drei. Nach Abspaltung seiner gemeinsamen U1-Ladung sechs bleibt
\begin{equation}
h(C)=\frac62-\frac{6^2}{2\cdot9}=1.
\end{equation}
Die C-Felder sind somit geladene Verzweigungsvertexfelder, keine gewöhnlichen neutralen Sechsfermionpolynome. Zusammen mit dem konjugierten $\Lambda_3$-Sektor ergeben sie genau die ursprüngliche lokale Erweiterung
\begin{equation}
V_{E_8}=V_{A_8}\oplus V_{A_8+\Lambda_3}\oplus V_{A_8+\Lambda_6},
\qquad248=80+84+84.
\end{equation}
Dieser bekannte mathematische Zusammenhang \cite{distler} ist hier an die tatsächlichen 24 Originalfelder angeschlossen. Die neun Hilfsfermionen sind eine Darstellungshilfe; ihre Anzahl behauptet weder neun Standardmodellteilchen noch eine neue räumliche Dimension.

\subsection{Die Sechser-Hyperkante ist ein elementarer Fermionring}
Für den bekannten Sechserzeugen gibt es in der Liste der zwölf elementaren Fermionfelder genau eine zulässige Wick-Paarung. Ihr Gesamtvorzeichen ist positiv. Sie liefert unabhängig von der Strom-Ward-Rekursion
\begin{equation}
F_6=\frac1{z_{12}z_{16}z_{23}z_{34}z_{45}z_{56}}=\frac1{576}
\end{equation}
an den bisherigen sechs Positionen. Alle 15 Paarungen ganzer Ströme bleiben null. Damit ist eine wichtige falsche Schlussfolgerung behoben: Nichtgaußische Stromantworten schließen eine zugrunde liegende Gaußdarstellung durch \emph{zusammengesetzte} Felder nicht aus. Sie schließen den Ersatz der Ströme durch lineare quasifreie Felder mit denselben Paarungen aus.

\picturetext{Sechs ganze Bauteile passen paarweise nicht zusammen. Jedes enthält aber zwei Anschlüsse. Über diese inneren Anschlüsse schließen sich alle sechs zu einem Ring. Wer jedes Bauteil durch einen einzigen Anschluss ersetzt, verliert den Ring und damit die berechnete gemeinsame Antwort.}

Das entspricht einem präzisen Hypergraphen: Die Sechserkante ist in den Stromvariablen irreduzibel; in einer reicheren, tatsächlich hergeleiteten Felddarstellung besitzt sie innere Struktur. Zuse- und Wolfram-Bilder können solche Darstellungen beschreiben. Der Wechsel der Darstellung muss die Antwort erhalten; er erzeugt keine neue Weltregel.

\subsection{Die volle C/X-Regel als geschlossenes Produkt}
Jedes Originalfeld ist ein phasenmarkierter Wurzelvertex $s_iV_{\alpha_i}$. Mit dem ursprünglichen bimultiplikativen E8-Kokzyklus $\varepsilon$ gilt bei verschiedenen endlichen Positionen
\begin{equation}\label{eq:latticeoracle}
F_n=\delta_{\sum_i\alpha_i,0}
\left(\prod_i s_i\right)
\left(\prod_{i<j}\varepsilon(\alpha_i,\alpha_j)\right)
\prod_{i<j}(z_i-z_j)^{(\alpha_i,\alpha_j)}.
\end{equation}
Die Potenzen sind ganzzahlig. Die Formel folgt direkt aus der Gittervertex-Konstruktion \cite{frenkelkac}; sie ist hier mit den Originalwurzeln, ihrer kompakten Adjungierung und ihren tatsächlichen Kokzyklusphasen implementiert. Bei neutraler Ladung enthält sie auch sämtliche geladenen Zwischenwege, die der Ward-Auswerter rekursiv durchläuft.

Ein neuer entscheidender Zeuge verlangt die Erweiterung wirklich:
\begin{equation}
C_0,C_1,C_2,X_{00},X_{11},X_{22},X_{33},X_{43}.
\end{equation}
Seine E8-Gesamtladung ist null. In der Hilfsfermiondarstellung beträgt die gemeinsame U1-Ladung jedoch 18. Ein gewöhnlicher Fermion-Vakuum-Wickwert wäre deshalb null; die volle ursprüngliche Quelle ergibt an $z=0,\ldots,7$
\begin{equation}
F_8=-\frac1{4032000}.
\end{equation}
Gitterprodukt und ursprüngliche Ward-Rekursion liefern denselben Wert. Die $\Z_3$-Erweiterung ist somit eine wirksame Komponente des Quellgesetzes, kein bloßes Etikett.

Für reine C/X-Wörter benötigt die Produktformel quadratisch viele exakte Paaroperationen. Sie vermeidet die wachsende Zahl rekursiver Ward-Verzweigungen. Die Bitkomplexität großer rationaler Zahlen kommt hinzu. Dies ist eine konkrete Vereinfachung des TFPT-Rechners, keine Behauptung einer neuen Erfindung der Gittervertexformel. Cartaneinsetzungen und beliebige Nachfahren sind nicht Eingaben dieser speziellen API; die ursprüngliche Ward-Regel bleibt für ihre entsprechenden Aufgaben vorhanden.

\subsection{Ein expliziter Universalraum für radiale Quellgeschichten}
Es sei
\begin{equation}
\ket{u}=\prod_{i=1}^m s_iV_{\alpha_i}(z_i)\Omega,
\qquad1>|z_1|>\cdots>|z_m|\ge0.
\end{equation}
Das ist eine radiale Feldgeschichte, keine bereits normierte Folge von Detektormessungen. Definiere
\begin{equation}
Q_u=\sum_i\alpha_i,\qquad p_{u,n}=\sum_i\alpha_i z_i^n\quad(n\ge1)
\end{equation}
und den Vorfaktor $A_u$ als das Produkt rechts in~\eqref{eq:latticeoracle} ohne Neutralitätsdelta. Die Heisenberg-Konstruktion gibt dann den \emph{vollständigen} Zustandsvektor
\begin{equation}\label{eq:fullhistory}
\ket{u}=A_u\exp\!\left(\sum_{n\ge1}\frac{p_{u,n}\cdot a_{-n}}{n}\right)\ket{Q_u}.
\end{equation}
Die Überlappung zweier solcher Geschichten ist daher
\begin{align}\label{eq:historyexact}
K(u,v)&=\delta_{Q_u,Q_v}\overline{A_u}A_v
\exp\!\left(\sum_{n\ge1}\frac{\overline{p_{u,n}}\cdot p_{v,n}}{n}\right)\nonumber\\
&=\delta_{Q_u,Q_v}\overline{A_u}A_v
\prod_{i,j}(1-\overline{z_{u,i}}z_{v,j})^{-(\alpha_{u,i},\alpha_{v,j})}.
\end{align}
Die Positivität folgt aus der tatsächlichen Hilbertraumdarstellung. Sie wurde zusätzlich an einer exakten Gram-Matrix und unabhängig über radial adjungierte Original-Ward-Wörter geprüft. Die implementierte API verwendet reelle rationale Positionen mit $|z|<1$, einschließlich negativer Positionen, in der angegebenen radialen Ordnung; die Formel erklärt den allgemeineren komplexen Radialbereich.

\key{Damit ist der Universalraum für diese Quellgeschichten konkret: Eine Geschichte hinterlässt ihre Gesamtladung, ihre Amplitude und ihre gesamte Modenfolge. Diese Daten bestimmen den unnormierten Zustandsvektor und damit alle späteren wohldefinierten linearen Quellantworten auf diesem Zustand.}

Die unendliche Modenfolge muss dabei nicht als unendliche Liste gespeichert werden. Ihre erzeugende Funktion ist rational:
\begin{equation}\label{eq:momentmemory}
P_u(t)=\sum_{n\ge1}p_{u,n}t^n
=\sum_i\alpha_i\frac{z_i t}{1-z_i t}.
\end{equation}
Eine endliche Geschichte besitzt also eine endliche exakte Beschreibung \emph{aller} dieser Nachfahrenmomente. Bei wachsender Geschichte kann die Zahl ihrer Pole wachsen; das ist kein Beweis eines universellen festen endlichen Gedächtnisrangs. Operationale Äquivalenz bei einer eingeschränkten Menge von Zukunftsmessungen kann außerdem gröber sein als die Gleichheit dieser unnormierten Vektoren.

Wird links ein neues Feld $s_\beta V_\beta(w)$ mit $1>|w|>|z_1|$ eingesetzt, lautet die genaue Fortschreibung
\begin{align}
Q'&=Q+\beta,&P'(t)&=P(t)+\beta\frac{wt}{1-wt},\\
A'&=A\,s_\beta\,\varepsilon(\beta,Q)
\prod_i(w-z_i)^{(\beta,\alpha_i)}.
\end{align}
Beim verwendeten bimultiplikativen Kokzyklus ist $\varepsilon(\beta,Q)=\prod_i\varepsilon(\beta,\alpha_i)$. Wenn eine Rekodierung diese Daten oder exakt denselben Zustandsvektor erhält, bleiben die nachfolgenden Quellantworten erhalten. Das schließt die Verbindung von geladenem Carry, positivem Kern und Nachfahrengedächtnis für diese konkrete Geschichtenklasse. Die Wahl physischer Messinstrumente und ihre Kausalnormierung bleibt eine zusätzliche Aufgabe.

\begin{center}
\begin{tikzpicture}[node distance=7mm]
\node[box,text width=13.2cm] (f) {\textbf{Dieselben ursprünglichen Felder}\quad $C_4+X_{20}$ mit Ladungen und Kokzyklus};
\node[box,below=of f,xshift=-3.45cm,text width=5.9cm] (w) {\textbf{Antworten}\quad Ward-Rekursion\quad $=$ Gitterprodukt für Wurzelwörter};
\node[box,below=of f,xshift=3.45cm,text width=5.9cm] (g) {\textbf{Geschichtszustand}\quad $Q$, $A$, rationale Funktion $P(t)$\quad alle Nachfahrenmomente};
\node[box,below=31mm of f,text width=13.2cm] (k) {\textbf{Positiver Universalraumkern}\quad $K(u,v)=\langle u|v\rangle$\quad exakte Fortschreibung und konforme Bewegung};
\node[box,below=of k,text width=13.2cm] (p) {\textbf{Noch zu identifizieren}\quad ursprüngliche physische Quelle und Raumzeitfelder\quad gleicher Zustand, gleiche Ladungen, gleiche Zeit};
\draw[arr](f)--(w);\draw[arr](f)--(g);\draw[arr](w)--(k);\draw[arr](g)--(k);\draw[dasharr](k)--(p);
\end{tikzpicture}
\end{center}

